Given a vector database $\mathcal X\subset\mathbb R^d$ and query $q$,
the task is to return $k$ approximate nearest neighbours under squared
Euclidean distance. The research problem is to construct and search
JHQ's compressed index within a single GPU's memory, and improve query
throughput at a given recall target.

Search first selects relevant IVF lists, screens their vectors using
primary codes, then uses residual codes to rescore a smaller shortlist.
The refinement budget controls how many candidates enter this last stage.
Ground truth (GT) gives exact neighbours in the original vector space;
recall is the fraction of these neighbours recovered.

The study retains JHQ's dimension-preserving orthogonal transform and
hierarchical codes (\cref{sec:2-3}), and investigates GPU construction,
data layout, search execution and refinement-budget selection.
Calibration compares smaller-budget outputs with a larger-budget
reference without GT, then reuses one chosen multiplier for a fixed
index and an assumed stable query distribution. GT evaluation checks
the resulting recall.

Evaluation uses one RTX 5090 and squared-L2 top-10. The main cross-system
comparison uses four datasets and batch 1,000; lookup-table, memory-layout
and budget tests cover six datasets. A small-batch study adds batches
8/32 (\cref{sec:7-smallbatch}). The index and core search computation
remain on the GPU; preparing and transferring queries within the search
call are included in its timing. Dynamic updates, host paging and
multi-GPU execution are outside the implementation; billion-scale
search is not evaluated.
